% ----------------------------------------------------------------------
\section{Introduction}
% ----------------------------------------------------------------------

% ----------------------------------------------------------------------
\begin{frame}
  \frametitle{Différentes classes de problèmes}

  \begin{itemize}
  \item Optimisation combinatoire\\
    \emph{branch and bound}, \emph{(meta)-heuristiques} 
  \item \alert{Optimisation continue}
    \begin{itemize}
    \item sans contrainte \\
      \emph{calcul analytique}, \emph{descente de gradient}
    \item avec contraintes  
      \begin{itemize}
        \item non-linéaires
        \item linéaires \\
        \emph{simplexe}
      \end{itemize}
    \end{itemize}
  \end{itemize}
  
\end{frame}

% ----------------------------------------------------------------------
\begin{frame}
  \frametitle{Formulation mathématique}

  \[
  \text{(P)} \left\{
  \begin{array}{c}
    \text{trouver} \ x \ \text{qui minimise} \ f(x) \ \text{tel que :} \\
    g_i(x) \leq 0, \ i \in I = \{1, \dots, m\} \\
    x \in S \subseteq \R^n
  \end{array}
  \right.
  \]

  \only<1>{
    \begin{block}{Vocabulaire}
    \begin{itemize}
    \item fonction objectif : $f : \R^n \rightarrow \R$
    \item inconnues : composantes du vecteur $x := (x_1, \dots, x_n) \in \R^n$
    \item contraintes : $g_i(x) \leq 0$ et $x \in S$
    \item une solution (candidate) : $x$ satisfaisant les contraintes
    \item solution optimale/optimum global : une solution minimisant $f$
    \end{itemize}
    %distinction des contraintes en deux:
    %- generalisation a plusieurs classes de pb
    %- traites differement dans les algos
    \end{block}
  }
  \only<2>{
    \begin{block}{Remarques}
    \begin{itemize}
    \item max $f'$ $\Leftrightarrow$ min $-f'$
    \item nombre d'inconnues $n \geq 1$ 
    \item nombre de contraintes $m \geq 0$
    \item contraintes d'égalité $g'(x) = 0$ $\Leftrightarrow$ $g'(x) \leq 0$
      et $-g'(x) \leq 0$
    \item généralement pas d'inégalités strictes
    \end{itemize}
    \end{block}
    %existence de solution (th. de Weierstrass)
    %si f est une fonction reelle continue sur K \subset R^n ferme et borne, alors il existe
    %une solution optimale x \in K qui min f(x)
    %existence aussi si f est une fonction reelle continue et possède la propriété de croissance à l'infini:
    %f(x) -> +inf qs x -> +inf
    %inégalités strictes => ensemble de solutions non ferme et borne, donc possiblement pas de solutions. 
  }
  \only<3>{
    \begin{block}{Résoudre (P) peut être difficile...}
    \begin{itemize}
    \item ensemble des solutions vide, très grand, voire infini
    \item problème non borné 
    \item $n, m$ grands
    \item $f$ et $g_i$ quelconques, $f$ non analytiquement connu %ni differentiable, ni convexe, ni lineaire
    %%\item $f$ peut être (quasiment) plat autour du min
      %si plat, plusieurs solutions optimales, si quasiment, instabilités numériques 
    \item cas discret ($S = \Z^n$), NP-complet 
    \end{itemize}
    \end{block}
  }

  \only<4->{
    \begin{block}{Résoudre (P) peut être facile...}
    \begin{itemize}
    \item cas continu
    \item problème borné
    \item une solution
    \item $f$ connu < différentiable < convexe
    \item aucune contrainte ou des contraintes linéaires
    \end{itemize}
    \end{block}
  }
\end{frame}
